% ============================================================
% Capitolo 3 — Design del sistema (traduzione italiana di
% chapters/ch03_system_design.tex; label invariati)
% ============================================================
\chapter{Design del sistema}
\label{ch:system}

Questo capitolo descrive la piattaforma sperimentale: un ambiente di
apprendimento per rinforzo multi-agente costruito sul simulatore CARLA,
in cui tre veicoli e tre pedoni sono addestrati congiuntamente con MAPPO
sotto due regimi di addestramento alternativi --- curriculum e
mixed-batch --- che condividono ogni altro componente dello stack.
L'obiettivo di design è che il regime di addestramento sia l'\emph{unica}
variabile indipendente dell'esperimento principale; tutto ciò che è
documentato in questo capitolo è comune a entrambi i bracci.

% ------------------------------------------------------------
\section{Panoramica dell'architettura}
\label{sec:architecture}

La piattaforma consiste di quattro strati (\cref{fig:architecture}): il
server CARLA, un ambiente multi-agente custom, lo stack di addestramento
RLlib e una pipeline di logging/valutazione.

\paragraph{Simulatore.}
CARLA 0.9.16~\cite{dosovitskiy2017carla}, costruito su Unreal Engine~4,
gira come server autonomo. L'ambiente lo pilota in \emph{modalità
sincrona} con un passo di fisica fisso di \SI{0.05}{s} (\SI{20}{Hz}): il
simulatore avanza solo quando l'ambiente emette un tick, il che rende la
raccolta di esperienza deterministica rispetto allo scheduling e
indipendente dalla velocità di orologio. Il traffico scriptato (non
apprendente) --- 15 veicoli NPC e 30 pedoni NPC in ogni livello di
addestramento --- è gestito dal Traffic Manager di CARLA, anch'esso in
modalità sincrona. L'addestramento usa la mappa urbana \emph{Town03};
\emph{Town05} è riservata esclusivamente al test di generalizzazione e
non viene mai vista durante l'addestramento.

\paragraph{Ambiente multi-agente.}
Un ambiente custom (\texttt{CarlaMultiAgentEnv}) espone la simulazione a
RLlib come interfaccia multi-agente con sei agenti apprendenti: tre
veicoli e tre pedoni (walker). A ogni reset l'ambiente spawna gli
agenti, assegna a ciascuno un percorso individuale (\cref{sec:routes}) e
configura lo scenario secondo il livello di difficoltà selezionato dal
manager curriculum/batch (\cref{sec:regimes}). Un episodio dura al più
1{.}000 passi di controllo (\SI{50}{s} di tempo simulato a \SI{20}{Hz}).

\paragraph{Stack di apprendimento.}
L'addestramento usa MAPPO~\cite{yu2022surprising} come implementato su
Ray/RLlib~2.10.0 con PyTorch~2.7 (Python 3.11.9), nel paradigma di
\emph{addestramento centralizzato con esecuzione decentralizzata}
(CTDE). Vengono apprese due policy: una \emph{policy veicolo} condivisa
dai tre veicoli e una \emph{policy pedone} condivisa dai tre pedoni
(condivisione di parametri entro ciascuna classe, nessuna condivisione
tra classi, poiché le due classi hanno spazi di osservazione e azione
diversi). Ogni actor è un percettrone multistrato con due strati nascosti
da 256 unità. Un critic centralizzato riceve un'osservazione globale a
225 dimensioni che concatena informazione per agente
(\cref{sec:obs-actions}); l'encoder del critic è un MLP di default, con
una variante GNN usata come asse sperimentale secondario
(\cref{sec:critic-variants}).

\paragraph{Logging e valutazione.}
Alla fine di ogni episodio il processo di addestramento accoda un record
JSON \emph{per agente} (sei per episodio) a un file
\texttt{episodes.jsonl}: ragione di terminazione, completamento del
percorso, efficienza di percorso, velocità e diagnostica di generazione
del percorso. Tutti i risultati del \cref{ch:results} sono calcolati da
questi record --- o dal flusso di addestramento (metriche cumulative,
on-policy) o dai record prodotti da una pipeline di valutazione
separata, che ricarica un checkpoint addestrato ed esegue quattro
scenari fissi (easy, medium, hard su Town03; test su Town05) da 100
episodi ciascuno (\cref{sec:critic-variants}).

\begin{figure}[t]
  \centering
  \begin{tikzpicture}[
      node distance=7mm and 12mm,
      box/.style={draw, rounded corners=2pt, align=center,
                  minimum width=52mm, minimum height=11mm, font=\small},
      lbl/.style={font=\scriptsize, midway},
      >={Stealth[length=2.5mm]}]
    \node[box] (carla) {Server CARLA 0.9.16\\ \scriptsize Town03 (train) / Town05 (test), TM};
    \node[box, below=of carla] (env) {\texttt{CarlaMultiAgentEnv}\\ \scriptsize 3 veicoli + 3 walker, route planner};
    \node[box, right=of env] (mgr) {Manager curriculum\\ / batch \scriptsize (livello per episodio)};
    \node[box, below=of env] (rllib) {Ray/RLlib 2.10 --- MAPPO (CTDE)\\ \scriptsize policy veicolo, policy pedone, critic centralizzato};
    \node[box, below=of rllib] (log) {\texttt{episodes.jsonl}\\ \scriptsize 6 record/episodio};
    \node[box, right=of log] (eval) {Pipeline di valutazione\\ \scriptsize $4$ scenari $\times$ 100 ep};
    \draw[->] (env) -- node[lbl, right] {tick / controllo} (carla);
    \draw[<-] ([xshift=-8mm]env.north) -- node[lbl, left] {sensori, stato} ([xshift=-8mm]carla.south);
    \draw[->] (mgr) -- node[lbl, above] {livello} (env);
    \draw[->] (env) -- node[lbl, right] {oss., reward} (rllib);
    \draw[<-] ([xshift=-8mm]env.south) -- node[lbl, left] {azioni} ([xshift=-8mm]rllib.north);
    \draw[->] (rllib) -- (log);
    \draw[->] (rllib) -| node[lbl, above, pos=0.25] {checkpoint} (eval);
  \end{tikzpicture}
  \caption{Architettura a componenti. Il manager curriculum/batch è
  l'unico componente che differisce tra i due bracci sperimentali, e
  solo nella regola di selezione del livello (\cref{sec:regimes}).}
  \label{fig:architecture}
\end{figure}

% ------------------------------------------------------------
\section{Un ambiente multi-agente a mondo condiviso}
\label{sec:shared-world}

\subsection{Copertura parallela dei task}
\label{sec:parallel-coverage}

I sei agenti apprendenti non cooperano verso un obiettivo congiunto: a
ogni agente è assegnato il \emph{proprio} percorso e ciascuno è valutato
individualmente sul suo completamento. Poiché tutti e sei i percorsi
sono collocati nello stesso mondo vivo, un singolo episodio esercita
simultaneamente sei micro-task diversi --- geometrie stradali diverse,
interazioni diverse con il traffico scriptato e, soprattutto,
interazioni reciproche tra agenti apprendenti (per es.\ un veicolo
apprendente che dà la precedenza a un pedone apprendente). Chiamiamo
questo design \emph{copertura parallela dei task} (parallel task
coverage): per episodio di orologio, il sistema raccoglie esperienza su
sei istanze di task anziché una.

Questo inquadramento porta con sé due avvertenze che l'analisi deve
rispettare. Primo, i sei campioni per episodio sono \emph{correlati} ---
gli agenti condividono lo stato del mondo, quindi un ingorgo o un
attraversamento bloccato influenza più record alla volta. Secondo, dal
punto di vista di ogni singolo agente l'ambiente è
\emph{non-stazionario} durante l'addestramento, perché le policy degli
altri cinque agenti cambiano nel tempo; è la motivazione standard del
critic CTDE (\cref{ch:background}). La copertura parallela dei task è
una proprietà della piattaforma, presente in modo identico in entrambi i
regimi di addestramento: è parte della cornice sperimentale, non una
variabile sperimentale.

\subsection{Percorsi per agente e seeding}
\label{sec:seeding}

A ogni reset, ciascun agente riceve un percorso indipendente la cui
lunghezza target è fissata dal livello di difficoltà attivo
(30/60/\SI{100}{m}; \cref{sec:routes}). La generazione dei percorsi è
seedata deterministicamente per agente con una sequenza di seed
costruita da
$\langle\text{seed traffico},\ \text{contatore di reset},\ \mathrm{crc32}(\text{id
agente})\rangle$, cosicché (i) due run di addestramento lanciate con lo
stesso seed di esperimento sullo stesso codice vedono la stessa sequenza
di estrazioni di generazione dei percorsi, e (ii) le estrazioni sono
disaccoppiate tra agenti. Questo rende i confronti A/B appaiati
accoppiati per percorso per costruzione, una proprietà su cui il
protocollo sperimentale fa affidamento (\cref{ch:methodology}).

\subsection{Tassonomia delle terminazioni e record per agente}
\label{sec:termination}

Ogni agente termina il proprio episodio in modo indipendente con
esattamente una di sei ragioni, classificate da una regola indipendente
dal simulatore (\texttt{episode\_classification.py}):

\begin{itemize}
  \item \texttt{route\_complete} --- l'agente ha raggiunto la fine del
        suo percorso. È l'\emph{unico} esito contato come successo.
  \item \texttt{route\_short} --- l'agente ha raggiunto la fine di un
        percorso la cui lunghezza ottima è risultata più corta del
        target di livello oltre un fattore di tolleranza (fallback
        degenere della generazione dei percorsi). Raggiungerla
        \emph{non} è contato come successo: questa guardia previene
        l'inflazione del tasso di successo tramite percorsi banalmente
        corti.
  \item \texttt{collision} --- contatto con qualunque attore o geometria
        statica.
  \item \texttt{offroad} --- il veicolo si è spostato più di \SI{5}{m}
        dal centro della corsia di guida più vicina (solo veicoli).
  \item \texttt{stuck} --- mancanza prolungata di progresso lungo il
        percorso.
  \item \texttt{timeout} --- è scaduto prima il tetto di 1{.}000 passi
        dell'episodio.
\end{itemize}

Ogni episodio produce quindi esattamente sei record a livello di agente
(tre veicolo, tre pedone) in \texttt{episodes.jsonl}. Tutte le metriche
aggregate di questa tesi sono calcolate a livello di agente su questi
record, sempre riportate separatamente per veicoli e pedoni; le regole
di misura, inclusa la deduplica e il criterio di successo rigoroso, sono
specificate nel \cref{ch:methodology}.

% ------------------------------------------------------------
% Fonti (verificate 2026-07-14): carla_multi_agent_env.py
%   _get_vehicle_obs L1291-1381, _get_pedestrian_obs L1383-1445,
%   _fill_nearby L1740-1764, costanti L84-88; layout obs globale
%   centralized_critic.py (compute_global_obs_dim_with_mask, L115).
\section{Spazi di osservazione e azione}
\label{sec:obs-actions}

Ogni agente riceve un vettore di osservazione compatto e specifico per
classe, con tutte le componenti normalizzate in $[-1, 1]$ e sanificate
contro valori non finiti prima di essere restituite all'apprendista. I
veicoli osservano 47 dimensioni e i pedoni 26; la \cref{tab:obs} elenca
entrambi i layout. Tre scelte di design meritano un commento.

\paragraph{Feature consapevoli del percorso.}
Oltre alle classiche feature di cinematica dell'ego e dei vicini più
prossimi, entrambe le classi osservano una breve \emph{anteprima} del
proprio percorso (offset relativi dei waypoint successivi nel sistema di
riferimento dell'ego) insieme all'errore di heading e, per i veicoli,
all'angolo della curva imminente e all'errore laterale con segno. Queste
feature danno alla policy la geometria locale del task assegnato senza
rivelare nulla degli obiettivi degli altri agenti.

\paragraph{Canale di anteprima dell'hazard.}
Una routine condivisa stima, lungo il corridoio frontale dell'agente, un
rischio di time-to-collision e una frazione di occupazione del
corridoio, calcolati separatamente contro i veicoli (corridoio largo
\SI{3.5}{m}, \SI{40}{m} in avanti, orizzonte \SI{5}{s}) e contro i
pedoni (\SI{3.0}{m}, \SI{30}{m}, \SI{4}{s}); i pedoni osservano solo il
canale veicoli (\SI{4.0}{m}, \SI{30}{m}, \SI{4}{s}). Gli stessi quattro
valori di rischio lato veicolo sono riusati dalla funzione di reward
come gate di sicurezza (\cref{sec:rewards}), così che la quantità che
modula la reward sia sempre visibile alla policy.

\paragraph{Feature di consistenza markoviana.}
Le ultime tre dimensioni del veicolo espongono stato interno
dell'ambiente da cui la reward dipende: il numero di passi senza
progresso sui waypoint (normalizzato per 300), lo stato della penalità
anti-loop e il tempo normalizzato rimanente prima del troncamento
dell'episodio. Senza di esse, le penalità della \cref{sec:rewards} e
l'orizzonte fisso agirebbero su informazione che l'agente non può
osservare --- un difetto di aliasing di stato identificato e corretto
durante lo sviluppo (\cref{ch:methodology}).

\begin{table}[t]
  \centering\small
  \caption{Layout delle osservazioni. Le barre indicano costanti di
  normalizzazione; le feature dei vicini sono espresse nel sistema di
  riferimento dell'ego.}
  \label{tab:obs}
  \begin{tabular}{@{}llp{8.0cm}@{}}
    \toprule
    Dim. & Gruppo & Contenuto \\
    \midrule
    \multicolumn{3}{@{}l}{\emph{Veicolo (47D)}} \\
    0--8   & Cinematica ego   & velocità (/30), accelerazione (/10), modulo velocità (/120 km/h), heading planare \\
    9--12  & Tracking percorso & distanza (/\SI{50}{m}) e bearing (/$\pi$) al waypoint successivo; offset laterale in corsia con segno (larghezze di corsia); progresso continuo sul percorso \\
    13--14 & Semaforo         & flag di presenza; stato (rosso 1.0, giallo 0.5, verde 0.0) \\
    15--23 & Veicoli vicini   & 3 più prossimi entro \SI{50}{m}: offset longitudinale e laterale (/50), velocità relativa (/30) \\
    24     & Memoria controllo & comando di sterzo precedente \\
    25--30 & Anteprima percorso & waypoint $+1/+3/+5$: offset longitudinale e laterale (/50) \\
    31--33 & Geometria percorso & errore di heading; angolo curva imminente; errore laterale con segno rispetto al percorso \\
    34--39 & Pedoni vicini    & 2 più prossimi entro \SI{50}{m}, stesse 3 feature \\
    40--43 & Anteprima hazard & rischio time-to-collision e occupazione del corridoio, vs.\ veicoli e vs.\ pedoni \\
    44--46 & Progresso/tempo  & passi senza progresso (/300); flag penalità anti-loop; tempo rimanente \\
    \midrule
    \multicolumn{3}{@{}l}{\emph{Pedone (26D)}} \\
    0--2   & Posizione        & posizione nel mondo ($x, y$ /500; $z$ /50) \\
    3--6   & Cinematica       & velocità (/5), modulo velocità (/5 m/s) \\
    7--8   & Heading          & vettore frontale planare \\
    9--10  & Waypoint         & distanza (/\SI{100}{m}) e bearing al waypoint corrente \\
    11     & Superficie       & flag marciapiede \\
    12--17 & Veicoli vicini   & 2 più prossimi entro \SI{50}{m}: stesse 3 feature \\
    18     & Progresso        & frazione di completamento del percorso \\
    19--22 & Anteprima percorso & waypoint $+1/+3$: offset longitudinale e laterale (/100) \\
    23     & Geometria percorso & errore di heading \\
    24--25 & Anteprima hazard & rischio time-to-collision e occupazione del corridoio, vs.\ veicoli \\
    \bottomrule
  \end{tabular}
\end{table}

\paragraph{Spazi di azione.}
Entrambe le classi usano azioni continue bidimensionali. L'azione del
veicolo è $[a_{\mathrm{tb}}, a_{\mathrm{steer}}] \in [-1,1]^2$, dove
$a_{\mathrm{tb}}$ positivo mappa sull'acceleratore e negativo sul freno
(la retromarcia è disabilitata), e $a_{\mathrm{steer}}$ è il comando di
sterzo. L'azione del pedone è
$[a_{\mathrm{speed}}, a_{\Delta\mathrm{h}}]$ con
$a_{\mathrm{speed}} \in [0,1]$ che scala una velocità massima di
camminata di \SI{5.0}{m/s} e $a_{\Delta\mathrm{h}} \in [-1,1]$ che
comanda un cambio di heading fino a $\pm 45^{\circ}$ per passo.

\paragraph{Osservazione globale per il critic centralizzato.}
Durante il solo addestramento, il critic riceve una concatenazione a
slot fissi delle osservazioni di tutti gli agenti seguita da una
maschera di vitalità per agente:
$[\,v_0\,|\,v_1\,|\,v_2\,|\,p_0\,|\,p_1\,|\,p_2\,|\,\text{alive}_{6}\,]$,
cioè $3 \times 47 + 3 \times 26 + 6 = 225$ dimensioni. La maschera di
vitalità è necessaria perché gli agenti terminano in modo indipendente
(\cref{sec:termination}): gli slot degli agenti terminati vengono
mascherati anziché riusati silenziosamente. Gli actor non vedono mai
questo vettore, preservando l'esecuzione decentralizzata.

% ------------------------------------------------------------
% Fonti (verificate 2026-07-14): carla_multi_agent_env.py
%   _vehicle_reward L1866-1960 ("Vehicle Reward v5"),
%   _pedestrian_reward L1962-2010 ("Pedestrian Reward v5").
\section{Design della reward}
\label{sec:rewards}

Le reward sono per agente e identiche nei due regimi di addestramento.
La \cref{tab:rewards} elenca ogni componente con trigger e ampiezza;
questa sezione ne spiega la struttura. La forma finale è l'esito delle
iterazioni guidate da gate documentate nella \cref{sec:iterations};
nessuna pretesa di ottimalità è attribuita ai singoli coefficienti.

\paragraph{Nucleo di progresso.}
Per entrambe le classi il segnale dominante è un bonus sparso per
waypoint del percorso raggiunto ($+100$ per i veicoli, $+50$ per i
pedoni), completato da un termine di shaping denso proporzionale alla
riduzione per passo della distanza dal waypoint successivo. I waypoint
saltati tagliando il percorso sono penalizzati, quindi il progresso non
può essere manipolato con scorciatoie.

\paragraph{Termini di sicurezza.}
Una collisione costa ai veicoli $-500$, applicata una volta al passo del
primo contatto (l'episodio poi termina). Due ulteriori termini per i
veicoli penalizzano l'uscita dalla corsia (lineare oltre \SI{2}{m} dal
centro corsia; la terminazione dell'episodio avviene a \SI{5}{m},
\cref{sec:termination}) e il superamento dei 50~km/h. Le collisioni dei
pedoni costano $-50$, e i pedoni ricevono un piccolo incentivo per passo
a restare sui marciapiedi anziché sulle corsie di guida.

\paragraph{Shaping anti-stallo condizionato all'hazard (veicoli).}
Il blocco empiricamente più delicato contrasta l'\emph{equilibrio
difensivo} in cui un veicolo si ferma indefinitamente per evitare
qualunque rischio. Tutti i suoi termini sono sospesi mentre il veicolo
attende legittimamente a un semaforo rosso o giallo. Lo scalare
$\mathit{hazard} =
\max(\text{TTC}_{\mathrm{veh}}, \text{occ}_{\mathrm{veh}},
\text{TTC}_{\mathrm{ped}}, \text{occ}_{\mathrm{ped}})$ riassume il
canale di anteprima dell'hazard della \cref{sec:obs-actions}, e un
booleano $\mathit{safe\_to\_push} = (\mathit{hazard} < 0.85)$ fa da gate
ai termini di spinta: una penalità di inattività sotto 2.5~km/h; un
bonus di partenza proporzionale ad acceleratore, allineamento al
percorso e a un fattore di urgenza che cresce con il tempo trascorso
senza progresso; e uno shaping di velocità minima attorno a un target di
8~km/h, attivo solo quando il veicolo è allineato al suo percorso
(allineamento $> 0.25$). Sopra la soglia di hazard ogni incentivo di
spinta scompare, quindi la policy non è mai premiata per accelerare in
un corridoio rischioso. Indipendentemente dal gate, una rampa di
non-progresso (dopo 100 passi senza raggiungere un waypoint) e una
penalità anti-loop rendono lo stallo prolungato e il girare in tondo
strettamente non redditizi, e una coppia di termini sulla fluidità dello
sterzo scoraggia il controllo a strattoni.

\paragraph{Banda di passo dei pedoni.}
I pedoni guadagnano un bonus per passo mentre camminano entro una banda
di comfort di $[1.2, 2.6]$~m/s, con una penalità di inattività sotto
0.15~m/s e una penalità di corsa sopra 3.0~m/s. La banda è
intenzionalmente ampia: a 2.0~m/s un pedone copre \SI{100}{m} --- il
target del livello hard --- entro l'orizzonte di \SI{50}{s}. Il design
asimmetrico (solo bonus) di questa banda e la sua interazione con il
gate di hazard dei veicoli sono analizzati empiricamente nel
\cref{ch:results}.

\begin{table}[t]
  \centering\small
  \caption{Componenti della reward. Le righe veicolo ``gated'' sono
  sospese durante l'attesa a un semaforo rosso/giallo; le righe marcate
  $^{\dagger}$ richiedono inoltre $\mathit{safe\_to\_push}$
  ($\mathit{hazard} < 0.85$) e allineamento al percorso $> 0.25$.}
  \label{tab:rewards}
  \begin{tabular}{@{}p{3.4cm}p{5.6cm}p{4.6cm}@{}}
    \toprule
    Componente & Trigger & Valore \\
    \midrule
    \multicolumn{3}{@{}l}{\emph{Veicolo}} \\
    Waypoint raggiunto    & per waypoint avanzato & $+100$ \\
    Waypoint saltato      & per waypoint saltato & $-10$ \\
    Progresso denso       & variazione della distanza dal waypoint successivo & $+4.0$ per metro guadagnato \\
    Collisione            & primo contatto (terminale) & $-500$ \\
    Fuori corsia          & $>$ \SI{2}{m} dal centro corsia & $-0.5$ per metro \\
    Eccesso di velocità   & velocità $> 50$ km/h & $-0.10$ per km/h oltre \\
    Inattività (gated)    & velocità $< 2.5$ km/h & $-0.15$ \\
    Bonus di partenza (gated$^{\dagger}$) & acceleratore da fermo & $(0.20 + 0.30\,u)\cdot\text{thr}\cdot\text{align}$, freno simmetrico; $u = \min(\text{no-progress}/200, 1)$ \\
    Shaping velocità minima (gated$^{\dagger}$) & velocità sotto / sopra 8 km/h & $-0.04\,(8 - v)\cdot\text{align}$ / $+0.15\cdot\text{align}$ \\
    Rampa di non-progresso (gated) & $>$ 100 passi senza waypoint & $-0.004$ per passo, tetto $-1.0$ \\
    Sterzo fluido         & $|\Delta \text{steer}| < 0.1$ a $> 5$ km/h & $+0.1$ \\
    Strattone di sterzo   & $|\Delta \text{steer}| > 0.5$ & $-0.3$ \\
    Anti-loop             & rilevatore di loop attivo & $-1.0$ \\
    \midrule
    \multicolumn{3}{@{}l}{\emph{Pedone}} \\
    Waypoint raggiunto    & per waypoint avanzato & $+50$ \\
    Waypoint saltato      & per waypoint saltato & $-10$ \\
    Progresso denso       & variazione della distanza dal waypoint & $+2.0$ per metro guadagnato \\
    Collisione            & primo contatto (terminale) & $-50$ \\
    Marciapiede / carreggiata & per passo su marciapiede / corsia di guida & $+0.2$ / $-0.3$ \\
    Passo di camminata    & velocità $\in [1.2, 2.6]$ m/s & $+0.3$ \\
    Inattività            & velocità $< 0.15$ m/s & $-0.15$ \\
    Corsa                 & velocità $> 3.0$ m/s & $-0.7$ per m/s oltre \\
    Anti-loop             & rilevatore di loop attivo & $-1.0$ \\
    \bottomrule
  \end{tabular}
\end{table}

% ------------------------------------------------------------
% Fonti (verificate 2026-07-14): route_planner.py (min_route_ratio=0.5,
%   max_route_ratio=2.0 L145-146/206-207, max_cross=3 L261/329, rigetto
%   catene corte L368-369); levels.yaml (levels_path 30/60/100, test
%   Town05 80); etichette route_source carla_multi_agent_env.py L1040-1159.
\section{Generazione dei percorsi e livelli di difficoltà}
\label{sec:routes}

La difficoltà è definita dalla \emph{lunghezza target} dei percorsi per
agente, con tutti gli altri parametri di scenario tenuti costanti: i tre
livelli di addestramento (easy, medium, hard) collocano i target di
veicoli e pedoni a 30, 60 e \SI{100}{m} rispettivamente, sempre su
Town03 e sempre con lo stesso traffico scriptato (15 veicoli NPC, 30
pedoni NPC). Un quarto livello di sola valutazione (\emph{test}) usa
Town05 con target di \SI{80}{m} e gli stessi conteggi di traffico. La
configurazione definisce anche famiglie di livelli in cui la difficoltà
varia per densità di traffico anziché per lunghezza del percorso; in
questa tesi non sono usate (\cref{sec:campaign}).

\paragraph{Percorsi dei veicoli.}
I percorsi dei veicoli sono pianificati sulla topologia stradale con il
global route planner di CARLA: le destinazioni candidate sono campionate
attorno alla distanza target e una candidata è accettata solo se la
lunghezza del suo cammino ottimo ricade in $[0.5, 2.0] \times$ il
target, con fino a 32 tentativi per reset. Se nessuna candidata
sopravvive, l'ambiente ripiega su un generatore legacy a catena di
waypoint. Il generatore applicato è registrato per record
(\texttt{route\_source} $\in$ \{\texttt{grp},
\texttt{legacy\_fallback}\}), quindi la quota di percorsi di fallback è
sempre misurabile.

\paragraph{Percorsi dei pedoni.}
I percorsi dei pedoni sono catene di segmenti di marciapiede connesse
attraverso attraversamenti pedonali. Due vincoli le plasmano: le catene
la cui lunghezza cumulativa scende sotto $0.5 \times$ il target vengono
rigettate (la stessa guardia inferiore che sostiene la retrocessione
\texttt{route\_short} della \cref{sec:termination}), e il numero di
attraversamenti per percorso è limitato a tre --- un tetto
all'esposizione dei pedoni sulla carreggiata il cui valore è stato
fissato empiricamente (\cref{sec:iterations}). Il generatore registra il
proprio esito per record (\texttt{sidewalk\_crosswalk},
\texttt{sidewalk\_distance}, \texttt{respawn}, \texttt{legacy\_chain}),
insieme a diagnostiche di fattibilità (lunghezza ottima, errore rispetto
al target, flag under-target e too-short, tentativi per candidata). La
fattibilità dei percorsi pedonali su Town03 è una limitazione
strutturale della piattaforma --- la topologia dei marciapiedi spesso
non può sostenere i target nominali medium/hard --- ed è analizzata
separatamente nella \cref{sec:ped-limits}, dove i limiti d'ambiente sono
esplicitamente distinti dai limiti di policy.

\todo{Fig.~3: percorsi di esempio per livello. Script generatore pronto
in \texttt{docs/thesis/scripts/make\_fig3\_route\_examples.py} (richiede
un server CARLA attivo; lo esegue l'utente); includere poi qui
\texttt{figures/route\_examples.pdf}.}

% ------------------------------------------------------------
% Fonti (verificate 2026-07-14): curriculum_batch.yaml (tutte le soglie);
%   curriculum_batch_manager.py (_balanced_policy_success_details L391-423:
%   minimo sulle policy; window_full L129/382; BatchLevelSampler L902-940).
\section{Regimi di addestramento curriculum e mixed-batch}
\label{sec:regimes}

Il regime di addestramento è la variabile indipendente dell'esperimento
principale. Entrambi i regimi sono implementati dallo stesso componente
manager, consultato una volta per reset di episodio per selezionare il
livello di difficoltà; condividono le definizioni di livello della
\cref{sec:routes}, il budget totale di $3 \times 10^6$ passi d'ambiente
e ogni altro componente descritto in questo capitolo. Differiscono
\emph{solo} nella regola di selezione del livello.

\subsection{Braccio curriculum: progressione condizionata alla competenza}
\label{sec:curriculum-arm}

L'addestramento parte con il solo livello easy sbloccato. I livelli
superiori si sbloccano attraverso un gate di competenza valutato su una
finestra scorrevole degli ultimi 50 episodi: il tasso di successo a
finestra è calcolato separatamente per la policy veicolo e la policy
pedone, e il gate vincola sul \emph{minimo} dei due --- la progressione
richiede che la classe più debole sia competente, non la media. Medium
si sblocca a un tasso di successo a finestra di almeno il 70\%, hard al
60\%; la finestra deve essere piena, e una quota minima del budget
totale deve già essere stata spesa ai livelli precedenti (20\% su easy
prima di medium; 28\% su medium prima di hard). Due decisioni di design
meritano giustificazione:

\begin{itemize}
  \item \textbf{A finestra, non cumulativo.} Il tasso di successo
        cumulativo è un integratore in ritardo trascinato verso il basso
        dai fallimenti di avvio a freddo: una policy può essere
        competente \emph{adesso} mentre la sua media dalla nascita è
        ancora ben sotto soglia. Il gating su finestra scorrevole misura
        la competenza corrente.
  \item \textbf{Cap di pressione.} Se una soglia non viene mai
        raggiunta, il livello si sblocca comunque forzatamente una volta
        consumata una quota fissa del budget totale (30\% per medium,
        70\% per hard). Questo garantisce che ogni livello sia visitato
        entro il budget, quindi il braccio curriculum non può mai
        stallare terminalmente su easy --- al peggio degrada verso
        un'esposizione mista ritardata.
\end{itemize}

Una volta sbloccati, i livelli sono campionati con pesi che normalizzano
a $0.30/0.35/0.35$ (easy/medium/hard) sotto vincoli di budget (easy al
più 30\% del budget una volta disponibile medium, medium al più 35\%,
hard almeno 35\%). I livelli appena sbloccati passano per una breve fase
di \emph{probation} con pressione di campionamento ridotta (hard è
temporaneamente de-enfatizzato a peso 0.85), rendendo la transizione
graduale anziché brusca.

\subsection{Braccio batch: alternanza stratificata}
\label{sec:batch-arm}

Il braccio mixed-batch è il \emph{controllo senza ordinamento}: tutti e
tre i livelli sono disponibili dal primo episodio. I livelli sono
campionati per mescolamento stratificato senza reinserimento con
finestra tre: ogni finestra consecutiva di tre episodi estratta da
un'istanza del campionatore è una permutazione uniformemente casuale di
\{easy, medium, hard\}, dando esposizione uniforme $1/3$ con uno
sbilanciamento massimo di un episodio per campionatore --- evitando sia
la deriva del campionamento i.i.d.\ sia qualunque ordinamento
deterministico fisso entro la finestra. Nei log della campagna
realizzata, dove gli episodi dei rollout worker paralleli si alternano,
le quote globali di livello restano entro $\approx 2$ punti percentuali
dall'uniforme e tutti e tre i livelli compaiono entro i primi venti
episodi (\cref{sec:integrity}).

\subsection{Che cosa isola il confronto}
\label{sec:isolation}

Sotto questa costruzione, i due bracci consumano lo stesso budget sulla
stessa distribuzione di task in attesa; ciò che differisce è
\emph{quando} ogni parte di quella distribuzione viene presentata e se
la presentazione è condizionata alla competenza misurata. Qualunque
differenza comportamentale tra i bracci è quindi attribuibile alla
politica di ordinamento, a meno della varianza da run a run --- che è
quantificata separatamente nel \cref{ch:results} e discussa nella
\cref{sec:threats}.

\begin{figure}[t]
  \centering
  \begin{tikzpicture}[
      state/.style={draw, rounded corners=2pt, align=center,
                    minimum width=34mm, minimum height=13mm, font=\small},
      gate/.style={font=\scriptsize, align=center},
      >={Stealth[length=2.5mm]}, node distance=17mm]
    \node[state] (s1) {solo easy\\ \scriptsize campionamento: easy 1.0};
    \node[state, right=of s1] (s2) {easy $+$ medium\\ \scriptsize probation, poi\\ \scriptsize $\propto 0.30/0.35$};
    \node[state, right=of s2] (s3) {easy $+$ medium $+$ hard\\ \scriptsize probation (hard $\times$0.85), poi\\ \scriptsize $\propto 0.30/0.35/0.35$};
    \draw[->] (s1.north east) to[bend left=28]
      node[gate, above] {SR fin.\ $\geq 0.70$\\ $\wedge$ quota easy $\geq 0.20$} (s2.north west);
    \draw[->, dashed] (s1.south east) to[bend right=28]
      node[gate, below] {cap: 30\% budget} (s2.south west);
    \draw[->] (s2.north east) to[bend left=28]
      node[gate, above] {SR fin.\ $\geq 0.60$\\ $\wedge$ quota med.\ $\geq 0.28$} (s3.north west);
    \draw[->, dashed] (s2.south east) to[bend right=28]
      node[gate, below] {cap: 70\% budget} (s3.south west);
  \end{tikzpicture}
  \caption{Progressione del curriculum. Archi pieni: gate di competenza
  sul tasso di successo a finestra (minimo sulle due policy, finestra di
  50 episodi, finestra piena) con soglie minime di quota di budget.
  Archi tratteggiati: cap di pressione che garantiscono la progressione
  entro il budget. Nella run MLP-curriculum della campagna (seed 999),
  medium si è sbloccato all'episodio 842 di 3{.}220 (26.1\%) e hard
  all'episodio 2{.}286 (71.0\%), per un mix finale di livelli di 32.5\%
  easy / 45.8\% medium / 21.6\% hard --- tempistiche coerenti con le
  repliche del trunk pre-campagna.}
  \label{fig:curriculum-fsm}
\end{figure}

% ------------------------------------------------------------
% Fonti (verificate 2026-07-14): centralized_critic.py (GraphConvLayer
%   L336-356, GNNCriticEncoder L414-497, critic MLP 225->256->256->1);
%   run_config.json delle run di campagna (tutti i valori
%   dell'ottimizzatore, byte-identici tra i bracci salvo mode/use_gnn).
\section{Varianti dell'encoder del critic e pipeline di addestramento}
\label{sec:critic-variants}

\subsection{Encoder del critic}
\label{sec:critic-encoders}

Nella configurazione di default il critic centralizzato è un MLP che
mappa l'osservazione globale a 225 dimensioni attraverso due strati
nascosti da 256 unità verso un valore scalare. L'asse sperimentale
secondario (\cref{sec:campaign}) sostituisce questo encoder con uno su
grafo, mentre \emph{entrambi gli actor restano gli identici MLP} della
\cref{sec:architecture}: qualunque differenza comportamentale tra le due
varianti è quindi attribuibile alla stima del valore durante
l'addestramento, non alla capacità della policy, e a tempo di esecuzione
le due varianti sono indistinguibili in architettura.

L'encoder GNN riseziona l'osservazione globale nei sei slot per agente,
proietta ciascuno attraverso uno strato lineare per classe in un
embedding di nodo a 64 dimensioni e applica due round di message passing
con aggregazione a media in stile
GraphSAGE~\cite{hamilton2017graphsage}: ogni nodo combina il proprio
embedding con la media degli embedding degli altri agenti \emph{vivi}
attraverso pesi appresi separati, seguita da un'attivazione ELU. Un
mean-pool mascherato sui nodi vivi produce l'input della testa di
valore. La maschera di vitalità dell'osservazione globale guida sia
l'aggregazione sia il pooling, così che gli agenti terminati escano dal
riassunto del critic invece di contribuire slot stantii. (La codebase
fornisce anche un encoder del critic basato su attenzione e la
normalizzazione del valore PopArt dietro flag di configurazione;
entrambi restano disabilitati in tutta questa tesi.)

\subsection{Impostazioni di ottimizzazione}
\label{sec:optimization}

Tutte le run della campagna condividono la stessa configurazione
dell'ottimizzatore, byte-identica tra i bracci salvo i flag di regime ed
encoder: learning rate $5\times10^{-4}$; sconto $\gamma = 0.997$; GAE
$\lambda = 0.95$; clip PPO 0.2 con penalità KL aggiuntiva (target 0.02);
gradient clipping a 0.5; coefficiente della value loss 0.05 con il clip
del valore di fatto disabilitato ($10^6$); batch di addestramento
8{.}000 passi, minibatch 256, 15 epoche SGD per batch, frammenti di
rollout da 200 passi. Il coefficiente di entropia segue uno schedule da
0.03 a 0.005 con la transizione ancorata all'83\% del budget totale
(passo 2.49M di 3M) --- espresso come frazione del budget così che run
diagnostiche brevi e run complete attraversino lo stesso profilo di
esplorazione. I file di configurazione completi sono riprodotti
nell'\cref{app:hyperparams}; la giustificazione empirica dei valori non
di default (in particolare $\gamma$, le impostazioni della value loss e
lo schedule di entropia) è parte del registro iterativo del
\cref{ch:methodology}.

\subsection{Pipeline di valutazione}
\label{sec:eval-pipeline}

I risultati sono misurati a tre livelli, tutti calcolati dai record a
livello di agente (\cref{sec:termination}):

\begin{enumerate}
  \item \textbf{Flusso di addestramento}: le metriche cumulative su
        tutti gli episodi di addestramento. Sono on-policy e, nel
        braccio curriculum, calcolate su un mix di livelli che cambia
        nel tempo --- caratterizzano il processo di addestramento, non
        l'abilità finale.
  \item \textbf{Ricalcolo statico}: un passaggio offline sui log di
        addestramento che ricalcola tutti gli aggregati da disco,
        deduplica i record e verifica l'integrità (esattamente sei
        record per episodio, nessun valore non finito).
  \item \textbf{Valutazione finale del checkpoint}: il checkpoint
        addestrato viene ricaricato ed eseguito su quattro scenari
        fissi da 100 episodi ciascuno --- easy, medium e hard su
        Town03, più il livello test su Town05 (\cref{sec:routes}),
        quest'ultimo a misurare la generalizzazione su strade mai viste
        in addestramento. Gli episodi sono valutati con la policy
        stocastica (campionando dalla gaussiana appresa) e re-seeding
        per episodio sia dello scenario sia del campionatore di azioni.
        Valutare la media deterministica della policy è stato
        deliberatamente rigettato: per una policy gaussiana continua
        collassa il comportamento su una singola traiettoria per
        scenario, un modo di fallimento analizzato nella
        \cref{sec:bugs}.
\end{enumerate}
